% Part: first-order-logic
% Chapter: semantics
% Section: intro-semantics

\documentclass[../../../include/open-logic-section]{subfiles}

\begin{document}

\olfileid[fr]{fol}{syn}{its}

\olsection{Introduction}

Attribuer un sens aux expressions relève de la sémantique. La notion
centrale est celle de satisfaction dans !!a{structure}. !!^a{structure}
donne un sens aux éléments de base du langage. Son !!{domain} est
un ensemble non vide d'objets : les quantificateurs portent sur ses
éléments, les !!{constant}s désignent des !!{element}s de ce domaine,
les !!{function}s sont interprétés par des fonctions d'une puissance
du domaine dans ce domaine, et les !!{predicate}s par des relations sur ce domaine.
Le !!{domain} et ces interprétations du vocabulaire de base constituent
!!a{structure}. Comme les !!{formula}s peuvent contenir des
!!{variable}s, nous devons aussi leur associer des !!{element}s du
!!{domain} : c'est une assignation de variables. La relation de
satisfaction réunit enfin ces données. Une !!{formula} peut être
satisfaite dans !!a{structure}~$\Struct{M}$ relativement à une
assignation de variables~$s$, ce que nous notons $\Sat{M}{!A}[s]$.
Cette relation est elle aussi définie par induction sur la structure
de~$!A$ : les tables de vérité des connecteurs logiques permettent,
par exemple, de définir la satisfaction de $(!A \land !B)$ à partir
de la satisfaction, ou de la non-satisfaction, de $!A$ et de~$!B$.
Nous verrons alors que l'assignation est sans incidence lorsque la
!!{formula}~$!A$ est !!a{sentence}, c'est-à-dire lorsqu'elle ne
possède aucune variable libre. Nous pourrons donc dire simplement
qu'un !!{sentence} est satisfait, ou non, dans !!a{structure}.

À partir de la relation de satisfaction $\Sat{M}{!A}$ pour les
!!{sentence}s, nous définirons les notions sémantiques fondamentales
de validité, de conséquence logique et de satisfaisabilité. Un
!!{sentence} est valide, $\Entails !A$, s'il est satisfait dans
toute !!{structure}. Il est conséquence d'un ensemble d'!!{sentence}s,
$\Gamma \Entails !A$, si toute !!{structure} qui satisfait tous les
!!{sentence}s de~$\Gamma$ satisfait également~$!A$. Enfin, un
ensemble d'!!{sentence}s est satisfaisable s'il existe une
!!{structure} qui les satisfait tous simultanément. Puisque les
!!{formula}s sont définies inductivement et que la satisfaction
l'est à son tour à partir de leur structure, nous pouvons démontrer
par induction des propriétés de cette sémantique et relier entre elles
les notions ainsi définies.

\end{document}
